\section{Proofs about algorithms (a few tips, including use of invariants)}\label{proofs: sec: algorithms}

\doHeader{Proofs about algorithms}

Proofs about \emph{algorithms} can be particularly challenging.
One reason is that during the execution of an algorithm its state changes.
Try to follow these tips when doing proofs about algorithms:

\begin{enumerate}
\item When you refer to the state of the algorithm
  (such as the value of some variable)
  make sure that it makes clear \emph{exactly what point in time}
  you have in mind.  

\item When you define an object in your proof (i.e., introduce a name for some quantity),
  by convention, the named object never changes after that --- it is static.
  Be aware of this convention and try to follow it.

  It is okay to violate this convention when you are referring in the proof
  to the name of variable in the algorithm.
  But (i) if you do this, make sure your wording signals it clearly.
  For example, say something like this:
  \begin{quote}
    1. Consider any execution of the algorithm. \\
    2. Consider the value of the algorithm's variable $x$ just before the first iteration.\\
    \vdots
  \end{quote}
  And (ii) whenever you refer to the value of $x$,
  make sure you make clear exactly \emph{when} during the execution you mean.

\item Learn to use \emph{invariants}.
  An invariant is a property that is maintained throughout the execution of the algorithm.
  Invariants are a basic proof technique for reasoning about iterative algorithms.
  Exercise~\ref{proofs: prob: elizabeth hides from john} in Section~\ref{proofs: sec: exercises} gives an example.
  We'll see more throughout these notes.

\item As a rule, reasoning about processes is generally harder than reasoning about static objects.
  So, when you can present your proof as reasoning only about static objects, do that.
  For example, in reasoning about a static object such as a set or a spanning tree,
  don't present the proof in a way that implicitly or explicitly introduces some new process
  (e.g.~one that builds the object step by step).
\end{enumerate}

Exercise~\ref{proofs: prob: elizabeth hides from john} in Section~\ref{proofs: sec: exercises}
illustrates the first three points above.

%%% Local Variables:
%%% mode: latex
%%% TeX-master: "long_form_proofs"
%%% End:
